\documentclass[10pt,letterpaper]{article}

\usepackage[T1]{fontenc}
\usepackage[utf8]{inputenc}
\usepackage[spanish,es-noquoting,es-noshorthands]{babel}
\usepackage{lmodern}
\usepackage[margin=1.9cm,top=1.5cm,bottom=1.6cm]{geometry}
\usepackage{booktabs}
\usepackage{array}
\usepackage{xcolor}
\usepackage{graphicx}
\usepackage{fancyvrb}

\definecolor{acc}{HTML}{1F4E79}
\definecolor{gris}{HTML}{5B6672}

\DeclareUnicodeCharacter{2192}{$\rightarrow$}

% Marca de algo pendiente (sale en rojo). Antes de entregar no debe quedar ningún uso.
\newcommand{\falta}[1]{{\color{red}\bfseries [FALTA: #1]}}

\setlength{\parskip}{2pt}
\setlength{\topsep}{2pt}
\setlength{\partopsep}{0pt}
\setlength{\parindent}{0pt}
\renewcommand{\arraystretch}{1.05}
\setlength{\emergencystretch}{4em}

% Secciones y pie de pagina sin paquetes externos: solo LaTeX base.
\makeatletter
\renewcommand\section{\@startsection{section}{1}{\z@}%
  {7pt \@plus 2pt \@minus 1pt}{3pt}%
  {\normalfont\normalsize\bfseries\color{acc}}}
\renewcommand\subsection{\@startsection{subsection}{2}{\z@}%
  {5pt \@plus 1pt}{2pt}%
  {\normalfont\small\bfseries\color{acc}}}
\renewcommand\@seccntformat[1]{\csname the#1\endcsname.\quad}
\renewcommand{\ps@plain}{%
  \renewcommand{\@oddhead}{}\renewcommand{\@evenhead}{}%
  \renewcommand{\@oddfoot}{\footnotesize\color{gris}%
    TI-4601 $\cdot$ Laboratorio 4 $\cdot$ Actas municipales en un grafo
    \hfill \thepage}%
  \renewcommand{\@evenfoot}{\@oddfoot}}
\makeatother
\pagestyle{plain}

\fvset{fontsize=\footnotesize,frame=leftline,framerule=0.6pt,rulecolor=\color{gris},xleftmargin=6pt}

\begin{document}

{\LARGE\bfseries\color{acc} Laboratorio 4 --- Actas municipales en un grafo}\\[2pt]
{\footnotesize\color{gris} TI-4601 Bases de Datos Avanzadas $\cdot$ Grupo 3 $\cdot$ Fabricio Mena $\cdot$ Alisson Redondo $\cdot$ Christopher Rodriguez}\\
{\footnotesize\color{gris} Repositorio: \texttt{https://github.com/fabrimena/Lab-4-Bases-de-Datos-Avanzadas} (notebook en \texttt{labs/lab4-neo4j/Lab4.ipynb})}

\vspace{4pt}\hrule\vspace{6pt}

\textbf{Entorno.} Laptop AMD Ryzen 7 5700U, 16 GB de RAM (4 CPU y 9,7 GiB para WSL 2), Windows 11 Home con
Ubuntu 24.04 en WSL 2 y Docker Desktop 29.5.2. Neo4j 5.26.32 Community, driver de Python 5.28.8 y
LlamaIndex 0.14.25. El notebook se ejecutó completo con \texttt{make lab4-ejecutar}.

\section{Delimitación (pregunta 1)}

\falta{Alisson: documentos del índice, actas cargadas y qué hacer con un PDF escaneado, con las cifras de la salida del paso 2}

\section{El grafo en el navegador y el cargo (pregunta 2)}

\falta{captura del paso 6.3: el camino de Gilberth a la 61-2026 y sus fragmentos, con el panel de propiedades de \texttt{PARTICIPA\_EN} abierto}

\texttt{Persona} representa a la persona y \texttt{PARTICIPA\_EN} representa su presencia en una sesión.
El cargo describe lo segundo: es lo que esa persona fue \emph{en esa acta}. La misma persona aparece en
varias actas; en nuestro grafo, 17 de las 18 personas están en las cuatro y una en tres.

Si el cargo fuera una propiedad de \texttt{Persona}, tendría un solo valor para todas sus sesiones.
Cuando cambie, por ejemplo porque se elige otra presidencia o porque en una sesión preside la
vicepresidenta, el \texttt{SET} de la carga sobrescribiría el valor y el grafo diría que la persona
siempre tuvo el último cargo cargado. Se pierde la historia y ya no se puede responder «¿quién presidió
la 58-2026?». En la relación, cada acta guarda su propio valor, y la consulta pregunta por la relación
y no por la persona. Es lo que comprueba la verificación del TODO 1:

\begin{Verbatim}
OK: el cargo vive en la relación: presidencia de 61-2026 = ['Gilberth González González']
\end{Verbatim}

En estas cuatro actas nadie cambió de cargo (lo revisamos con una consulta aparte): Gilberth González
tiene cuatro relaciones con \texttt{cargo = 'Presidente'}, una por acta. Con más actas sí habría
cambios, y el modelo ya los admite sin tocar el esquema. Lo mismo vale para \texttt{rol}: alguien puede
ser suplente en un periodo y propietario en otro.

\section{Consultas de uno y dos saltos (pregunta 3)}

\falta{Alisson: consultas de los TODO 3 y 4 con su resultado, el SQL equivalente y cuántos \texttt{JOIN} necesitaría un tercer salto}

\section{\texttt{CREATE} contra \texttt{MERGE} (pregunta 4)}

Estas son nuestras consultas de los TODO 1 y 2. Siguen el ejemplo de las actas: \texttt{MATCH} del acta,
\texttt{MERGE} por la clave única y \texttt{SET} del resto.
\begin{Verbatim}
// TODO 1                                     // TODO 2
MATCH (a:Acta {id: $id})                      MATCH (a:Acta {id: $acta})
MERGE (p:Persona {nombre: $nombre})           MERGE (f:Fragmento {id: $id})
MERGE (p)-[r:PARTICIPA_EN]->(a)               SET f.indice = $indice, f.texto = $texto,
SET r.rol = $rol, r.cargo = $cargo                f.embedding = $embedding
                                              MERGE (a)-[:CONTIENE]->(f)
\end{Verbatim}
Salida del notebook. Cada verificación carga dos veces y compara los conteos:
\begin{Verbatim}
OK: cargar dos veces no duplica nada: (18, 71) → (18, 71)
OK: (personas, relaciones) = (18, 71); se esperaba (18, 71)
OK: cargar dos veces no duplica fragmentos ni relaciones
OK: conteos {'fragmentos': 15, 'contiene': 15, 'sigue': 11};
    se esperaba {'fragmentos': 15, 'contiene': 15, 'sigue': 11}
OK: todos los fragmentos tienen su vector
\end{Verbatim}

Para ver qué pasa con \texttt{CREATE}, corrimos una copia del notebook hasta el paso 5.3 con dos
variantes del TODO 1. El notebook del repositorio no se tocó.

\textbf{a) \texttt{CREATE} para la persona y para la relación.}
\begin{Verbatim}
ConstraintError: {code: Neo.ClientError.Schema.ConstraintValidationFailed} {message: Node(0) already
exists with label `Persona` and property `nombre` = 'Gilberth González González'}
VERIFICACIÓN FALLÓ: la carga intentó crear de nuevo un nodo que ya existía. Use MERGE (no CREATE)
para la persona y para la relación.
conteo (personas, relaciones): (18, 18)
\end{Verbatim}
No llega a la segunda carga: falla en la primera, al pasar a la segunda acta. La 58-2026 crea sus 18
personas, pero el primer participante de la 59-2026 es Gilberth, que ya existe, y la restricción
\texttt{persona\_nombre} rechaza el segundo nodo. Como cada \texttt{correr} es su propia transacción, lo
anterior queda guardado y el grafo queda a medias: solo la 58-2026 tiene participantes. Sin la
restricción no habría error, pero quedarían 71 nodos \texttt{Persona} en lugar de 18, uno por
participación, y las cuatro sesiones de Gilberth no estarían unidas por el mismo nodo.

\textbf{b) \texttt{MERGE} para la persona y \texttt{CREATE} para la relación.}
\begin{Verbatim}
VERIFICACIÓN FALLÓ: cargar dos veces no duplica nada: (18, 71) → (18, 142)
\end{Verbatim}
Aquí Neo4j no da ningún error, porque la restricción es sobre \texttt{Persona.nombre} y no hay ninguna
sobre relaciones. La primera carga deja bien los conteos, pero la segunda duplica las 71 relaciones:
Gilberth queda con dos \texttt{PARTICIPA\_EN} hacia cada acta, las dos con \texttt{cargo = 'Presidente'},
y cualquier conteo de participaciones sale doble. Es el caso más peligroso porque falla en silencio; solo
lo detecta la verificación que carga dos veces.

Con \texttt{MERGE} en los dos lugares la carga es idempotente: la segunda vez encuentra los mismos nodos
y relaciones y solo vuelve a aplicar el \texttt{SET}. La restricción protege los nodos, pero las
relaciones solo las protege el \texttt{MERGE}.

\section{Solo vector contra vector + acta (pregunta 5)}

\falta{Christopher: salida del TODO 5, comparación de «solo vector» con «vector + acta» y el papel de los 128 tokens con la salida del paso 8.2}

\section{Lectura crítica del modelo (pregunta 6)}

\falta{Christopher: las tres preguntas del paso 9, si cada respuesta es correcta según el contexto recuperado y qué afirma el modelo que no está en ese contexto}

\end{document}
